%%%%%%%%%%%%%%%%%%%%%%%%%%%%%%%%%%%%%%%%%%%%%%%%%%%%%%%%%%%%%%%%%%%%%%%
%%  Farbpalette des RWTH Corporate Design
%%
%%  Quelle: offizielle CD-Farbdefinitionen der RWTH Aachen,
%%          gespiegelt unter github.com/ifs-rwth-aachen/RWTH-Colors
%%          (LaTeX/rwthcolors.tex und CSS/rwth-colors.css, gegengeprueft
%%           am 21.08.2026 - beide Quellen stimmen exakt ueberein).
%%
%%  Die RWTH definiert eine Hausfarbe (Blau) und elf Sekundaerfarben,
%%  jede in fuenf Abstufungen: 100 / 75 / 50 / 25 / 10 Prozent.
%%  Die Abstufungen sind KEINE rechnerischen Aufhellungen, sondern vom
%%  CD einzeln festgelegt - daher hier explizit ausgeschrieben und
%%  nicht ueber xcolor-Mischungen (blau!75) erzeugt.
%%
%%  Benennung: RWTH<Farbe><Stufe>, z. B. \color{RWTHBlau75}
%%  Zusaetzlich unten Kurz-Aliase fuer die im Text genutzten Farben.
%%
%%  Benoetigt: \usepackage{xcolor} VOR diesem \input.
%%%%%%%%%%%%%%%%%%%%%%%%%%%%%%%%%%%%%%%%%%%%%%%%%%%%%%%%%%%%%%%%%%%%%%%

%% Hausfarbe -----------------------------------------------------------
\definecolor{RWTHBlau100}{RGB}{0,84,159}       % 00549F  Logo-Primaerfarbe
\definecolor{RWTHBlau75}{RGB}{64,127,183}      % 407FB7
\definecolor{RWTHBlau50}{RGB}{142,186,229}     % 8EBAE5  Logo-Sekundaerfarbe
\definecolor{RWTHBlau25}{RGB}{199,221,242}     % C7DDF2
\definecolor{RWTHBlau10}{RGB}{232,241,250}     % E8F1FA

%% Sekundaerfarben -----------------------------------------------------
% Schwarz (die 75%-Stufe ist das CD-Grau)
\definecolor{RWTHSchwarz100}{RGB}{0,0,0}       % 000000
\definecolor{RWTHSchwarz75}{RGB}{100,101,103}  % 646567
\definecolor{RWTHSchwarz50}{RGB}{156,158,159}  % 9C9E9F
\definecolor{RWTHSchwarz25}{RGB}{207,209,210}  % CFD1D2
\definecolor{RWTHSchwarz10}{RGB}{236,237,237}  % ECEDED

% Magenta
\definecolor{RWTHMagenta100}{RGB}{227,0,102}   % E30066
\definecolor{RWTHMagenta75}{RGB}{233,96,136}   % E96088
\definecolor{RWTHMagenta50}{RGB}{241,158,177}  % F19EB1
\definecolor{RWTHMagenta25}{RGB}{249,210,218}  % F9D2DA
\definecolor{RWTHMagenta10}{RGB}{253,238,240}  % FDEEF0

% Gelb
\definecolor{RWTHGelb100}{RGB}{255,237,0}      % FFED00
\definecolor{RWTHGelb75}{RGB}{255,240,85}      % FFF055
\definecolor{RWTHGelb50}{RGB}{255,245,155}     % FFF59B
\definecolor{RWTHGelb25}{RGB}{255,250,209}     % FFFAD1
\definecolor{RWTHGelb10}{RGB}{255,253,238}     % FFFDEE

% Petrol
\definecolor{RWTHPetrol100}{RGB}{0,97,101}     % 006165
\definecolor{RWTHPetrol75}{RGB}{45,127,131}    % 2D7F83
\definecolor{RWTHPetrol50}{RGB}{125,164,167}   % 7DA4A7
\definecolor{RWTHPetrol25}{RGB}{191,208,209}   % BFD0D1
\definecolor{RWTHPetrol10}{RGB}{230,236,236}   % E6ECEC

% Tuerkis
\definecolor{RWTHTuerkis100}{RGB}{0,152,161}   % 0098A1
\definecolor{RWTHTuerkis75}{RGB}{0,177,183}    % 00B1B7
\definecolor{RWTHTuerkis50}{RGB}{137,204,207}  % 89CCCF
\definecolor{RWTHTuerkis25}{RGB}{202,231,231}  % CAE7E7
\definecolor{RWTHTuerkis10}{RGB}{235,246,246}  % EBF6F6

% Gruen
\definecolor{RWTHGruen100}{RGB}{87,171,39}     % 57AB27
\definecolor{RWTHGruen75}{RGB}{141,192,96}     % 8DC060
\definecolor{RWTHGruen50}{RGB}{184,214,152}    % B8D698
\definecolor{RWTHGruen25}{RGB}{221,235,206}    % DDEBCE
\definecolor{RWTHGruen10}{RGB}{242,247,236}    % F2F7EC

% Maigruen
\definecolor{RWTHMaigruen100}{RGB}{189,205,0}  % BDCD00
\definecolor{RWTHMaigruen75}{RGB}{208,217,92}  % D0D95C
\definecolor{RWTHMaigruen50}{RGB}{224,230,154} % E0E69A
\definecolor{RWTHMaigruen25}{RGB}{240,243,208} % F0F3D0
\definecolor{RWTHMaigruen10}{RGB}{249,250,237} % F9FAED

% Orange
\definecolor{RWTHOrange100}{RGB}{246,168,0}    % F6A800
\definecolor{RWTHOrange75}{RGB}{250,190,80}    % FABE50
\definecolor{RWTHOrange50}{RGB}{253,212,143}   % FDD48F
\definecolor{RWTHOrange25}{RGB}{254,234,201}   % FEEAC9
\definecolor{RWTHOrange10}{RGB}{255,247,234}   % FFF7EA

% Rot
\definecolor{RWTHRot100}{RGB}{204,7,30}        % CC071E
\definecolor{RWTHRot75}{RGB}{216,92,65}        % D85C41
\definecolor{RWTHRot50}{RGB}{230,150,121}      % E69679
\definecolor{RWTHRot25}{RGB}{243,205,187}      % F3CDBB
\definecolor{RWTHRot10}{RGB}{250,235,227}      % FAEBE3

% Bordeaux
\definecolor{RWTHBordeaux100}{RGB}{161,16,53}  % A11035
\definecolor{RWTHBordeaux75}{RGB}{182,82,86}   % B65256
\definecolor{RWTHBordeaux50}{RGB}{205,139,135} % CD8B87
\definecolor{RWTHBordeaux25}{RGB}{229,197,192} % E5C5C0
\definecolor{RWTHBordeaux10}{RGB}{245,232,229} % F5E8E5

% Violett
\definecolor{RWTHViolett100}{RGB}{97,33,88}    % 612158
\definecolor{RWTHViolett75}{RGB}{131,78,117}   % 834E75
\definecolor{RWTHViolett50}{RGB}{168,133,158}  % A8859E
\definecolor{RWTHViolett25}{RGB}{210,192,205}  % D2C0CD
\definecolor{RWTHViolett10}{RGB}{237,229,234}  % EDE5EA

% Lila
\definecolor{RWTHLila100}{RGB}{122,111,172}    % 7A6FAC
\definecolor{RWTHLila75}{RGB}{155,145,193}     % 9B91C1
\definecolor{RWTHLila50}{RGB}{188,181,215}     % BCB5D7
\definecolor{RWTHLila25}{RGB}{222,218,235}     % DEDAEB
\definecolor{RWTHLila10}{RGB}{242,240,247}     % F2F0F7

%%%%%%%%%%%%%%%%%%%%%%%%%%%%%%%%%%%%%%%%%%%%%%%%%%%%%%%%%%%%%%%%%%%%%%%
%%  Kurz-Aliase fuer die im Fliesstext und in den TikZ-Grafiken
%%  verwendeten Farben. Ohne Stufenangabe ist immer 100 % gemeint.
%%  Aendert man hier den Bezug, aendert sich die Grafik dokumentweit.
%%%%%%%%%%%%%%%%%%%%%%%%%%%%%%%%%%%%%%%%%%%%%%%%%%%%%%%%%%%%%%%%%%%%%%%
\colorlet{rwthblau}{RWTHBlau100}
\colorlet{rwthblau75}{RWTHBlau75}
\colorlet{rwthblau50}{RWTHBlau50}
\colorlet{rwthblau25}{RWTHBlau25}
\colorlet{rwthblau10}{RWTHBlau10}
\colorlet{rwthgruen}{RWTHGruen100}
\colorlet{rwthorange}{RWTHOrange100}
\colorlet{rwthbordeaux}{RWTHBordeaux100}
\colorlet{rwthgrau}{RWTHSchwarz75}   % das CD-Grau ist Schwarz 75 %
